\chapter{Specific requirements}\label{ch:requirements}
%======================================================================

\section{Scope and conventions}

This chapter defines specific requirements of the \ssat{} that are
conditions for the validation against the requirements baseline.

All the requirements are to be considered as ``essential'' unless
otherwise stated.

Baseline, or customer requirements identified in this document are
referenced as \textbf{BR\#} (Baseline Requirement).

%----------------------------------------------------------------------
\section{System functional requirements}\label{sc:functional}
%----------------------------------------------------------------------

\begin{requirement}{Ada and C support}
  The tool will accept source code written in Ada or C.
\end{requirement}

\begin{rationale}
  The GNAT Pro development environment supports these two programming
  languages.
\end{rationale}

\begin{requirement}{Per subprogram stack usage information}
  \ssat{} will be able to extract the maximum stack usage information
  for each individual subprogram. This number will be qualified by an
  attribute indicating whether the maximum for this subprogram is a
  fixed number that has been statically determined or there is an
  additional amount that can only be determined at run time.
\end{requirement}

\begin{rationale}
  This information is very valuable for understanding which are
  the subprograms that are consuming more stack, or those whose stack
  requirements cannot be statically determined.
\end{rationale}

\begin{requirement}{Per subprogram call chain information}
  \ssat{} will generate, for each subprogram, the list of subprograms
  that it may call.
\end{requirement}

\begin{rationale}
  This information is needed for building the complete call graph
  containing all the possible call paths that the application can
  follow.
\end{rationale}

\begin{requirement}{Call graph information}
  \ssat{} will be able to traverse all possible call paths and sort
  them according to their global stack usage.
\end{requirement}

\begin{rationale}
  The longest call path, in terms of stack usage, is the parameter
  that we need for dimensioning the stack. In addition, having access
  to the rest of the call paths is useful because the most stack
  consuming path may not be feasible (because of run-time constraints
  known by the user).
\end{rationale}

\begin{requirement}{Support for multiple entry points}
  \ssat{} will be able to use any number of entry points from which
  different call graphs are computed. Tasks will be automatically used
  by \ssat{} as entry points, and users may add any number of entry
  points.
\end{requirement}

\begin{rationale}
  Tasks have their own stack, and hence per task stack usage is needed
  for dimensioning each individual stack. The ability to compute stack
  requirements for user defined entry points is also useful for
  documenting the stack requirements of any piece of code.
\end{rationale}

\begin{requirement}{Detect targets for indirect calls}
  Indirect calls are those whose target address can only be determined
  at run time, including dispatching calls and dereferences of
  access-to-subprogram values. \ssat{} will generate a list of
  candidate target nodes for each indirect call.
\end{requirement}

\begin{rationale}
  A list of every potential target address for indirect calls is of
  great value for obtaining safe worst case values.
  % \ssat{} will use the list of access, address, and overloaded
  % values and compare them agaisnt the expected signature. This way
  % some pessimism is removed from the stack usage computation
\end{rationale}

\begin{requirement}{Bound the size of dynamic frames}
  Dynamic frames are those whose stack usage cannot be determined
  statically because it depends on run-time information. Users will be
  able to define the maximum expected bounds to these variable-size
  frames.
\end{requirement}

\begin{rationale}
  When \ssat{} is not able to statically determine the stack
  requirements for a certain frame then the safest approach is to
  allow the user to determine a maximum expected bound.
\end{rationale}

\begin{requirement}{Support for generics}
  \ssat{} will extract stack requirements for any instantiation of a
  generic package.
\end{requirement}

\begin{rationale}
  Generic packages are widely used and a very valuable feature of the
  Ada language.
\end{rationale}

\begin{requirement}{Support for exceptions}
  \ssat{} will compute stack requirements for exception handlers and
  will use these values for generating the different call paths.
\end{requirement}

\begin{rationale}
  Potential stack consumption made by exception handlers must always
  be taken into account for achieving accurate stack analysis. 
\end{rationale}

\begin{requirement}{Support for interrupts}
  \ssat{} will compute stack requirements for interrupt handlers
  (protected procedure handlers).
\end{requirement}

\begin{rationale}
  Potential stack consumption made by interrupt handlers must always
  be taken into account for achieving accurate stack
  analysis. This information is needed for dimensioning the interrupt
  stacks (if they use separate stacks) or the different task stacks.
\end{rationale}

\begin{requirement}{Define manually stack requirements for subprograms}
  \ssat{} will allow users to specify maximum stack requirements for
  any individual subprogram. This value will be used for computing
  global stack requirements.
\end{requirement}

\begin{rationale}
  This feature is very useful when we have third party object
  libraries for which we cannot use \ssat{} for obtaining their stack
  information.
\end{rationale}

\begin{requirement}{Modify manually the list of targets for indirect calls}
  \ssat{} will allow users to modify, for any indirect call, the list of
  possible subprograms that may be called. This information will be added to
  the list of target subprograms that is automatically computed by \ssat{}. 
\end{requirement}

\begin{rationale}
  The list of possible target subprograms that is automatically
  computed by \ssat{} is based on static information, and hence it may
  be very pessimistic. The user will be allowed to restrict the list
  possible subprograms that may be called from a certain indirect
  call.
\end{rationale}

\begin{requirement}{Bound cycles in the call graph}
  A cycle in the call graph refers to a call sequence that may reach a
  node already belonging to the call sequence. The user will be able
  to specify the maximum number of cycle iterations so that the stack
  analysis tool can generate a safe maximum for the stack usage.
\end{requirement}

\begin{rationale}
  In presence of control flow cycles involving stack consumption, the
  worst-case stack usage is a function of the maximum number of cycle
  iterations.
\end{rationale}

\begin{requirement}{Manual exclusion of subprograms from the analysis}
  Users will be able to define a set of subprograms that will not be
  analyzed by \ssat{}. 
\end{requirement}

\begin{rationale}
  There may be cases that \ssat{} provides very pessimistic results
  and users may prefer to take advantage of run-time limitations that
  are not known by \ssat{}.
\end{rationale}

%----------------------------------------------------------------------
\section{System interface requirements}\label{sc:interfaces}
%----------------------------------------------------------------------

\subsection{System level interfaces}\label{sc:system-interfaces}

\begin{requirement}{Development platforms}
  The \ssat{} will run on any development platform supported by AdaCore.
\end{requirement}

\begin{rationale}
  The tool will be tightly tied to the compiler itself.
\end{rationale}

\begin{requirement}{Integrate with GNAT's Project Manager}
  GNAT's Project Manager is a facility that allows you to manage
  complex builds involving a number of source files, directories, and
  compilation options for different system configurations.  
 \ssat{} will be usable on projects configured using the GNAT's
  Project Manager.
\end{requirement}

\begin{rationale}
  GNAT's Project Manager is the natural way of handling complex
  systems when using the GNAT Pro tools.
\end{rationale}

\begin{requirement}{Textual report}
  \ssat{} will generate textual files containing per subprogram
  information (stack requirements and call chain information), as well
  as call graph information containing global stack usage.
\end{requirement}

\begin{rationale}
  Textual are very useful for being parsed by external analysis
  scripts and for being added to any generated documentation.
\end{rationale}

\begin{requirement}{Graphic interface}
  A graphic interface based on current GPS technology will be used on
  the development platform to access the compilation and the debugging
  tools.
\end{requirement}

\begin{rationale}
  GPS is the graphical IDE supported by AdaCore and it will allow a
  simple and efficient interaction between the source code and the
  results of the \ssat{}.
\end{rationale}

%----------------------------------------------------------------------
\subsection{System level data interfaces}\label{sc:data-interfaces}
%----------------------------------------------------------------------

\begin{requirement}{Source files}
  Ada and C source programs will be represented in standard text files,
  using Latin-1 coding \refcite{GNATUG}. Each file will contain a single
  Ada or C compilation unit, including any pragmas associated with the
  unit.
  % The default file name  determined by the name of the unit that the
  % file contains. The name is formed by taking the full expanded name
  % of the unit and replacing the separating dots with hyphens and using
  % lowercase for all letters.
\end{requirement}

\begin{rationale}
  Latin-1 is an 8-bit code that includes the familiar 7-bit ASCII set,
  plus additional characters used for representing foreign languages.
  % The default file extensions will be `.ads' for a spec and `.adb' for
  % a body. However, arbitrary file names may also be used.
\end{rationale}

\begin{requirement}{Generated textual report files}
  \ssat{} will generate standard text files, using Latin-1 coding
  \refcite{GNATUG}.
\end{requirement}

\begin{rationale}
  Latin-1 is an 8-bit code that includes the familiar 7-bit ASCII set,
  plus additional characters used for representing foreign languages.
\end{rationale}

\begin{requirement}{Generated graphical output}
  \ssat{} will generate an annotated call graph containing the
  different subprograms (nodes), their stack requirements
  (annotations), and possible call chains (edges).
  The VCG \refcite{VCG} graph description language, already used in
  GCC for other graph outputs, will be used for this purpose.
\end{requirement}

\begin{rationale}
  This standalone graph will be very useful both for understanding the
  program and for documentation purposes.
\end{rationale}

%----------------------------------------------------------------------
\section{Software configuration management}\label{sc:configuration}
%----------------------------------------------------------------------

\begin{requirement}{Identification of reference versions}
  All such reference versions will be identified by a specific, unique
  identifier.
\end{requirement}

\begin{rationale}
  AdaCore implements a configuration management policy (currently
  based on CVS), and identifies with tags key points in the software
  development process. Examples of reference versions include all GNAT
  Pro official releases.
\end{rationale}

\begin{requirement}{Reconstruction of past states of the software}
  It will be possible to reconstruct any reference version of the
  software.
\end{requirement}

\begin{rationale}
  Should the need of reconstructing a reference version of the
  compiler arise, the following procedure would fulfil that need:
  \begin{itemize}
    \item Get the tag corresponding to the reference version
    \item Get the sources corresponding to this tag
    \item Get the software construction tools corresponding to this
      tag
    \item Build the sources using the tools
  \end{itemize}
  AdaCore has standard procedures to carry out all these tests and uses
  them on a day-to-day basis.
\end{rationale}

\begin{requirement}{Tracking changes between two different states of
  the software}
  It will be possible to analyze differences between two different
  arbitrary versions of the software.
\end{requirement}

\begin{rationale}
  AdaCore also has standard procedures for doing so. CVS has a
  low-level interface that enables comparison between different
  versions of any file.
  However this mechanism is limited, because it is difficult to have a
  global view of all the changes related to a single technical
  change. For instance, implementing one feature may require changes
  in many files.
  To workaround this, AdaCore's company policy is to assign a tracking
  number to every feature to be included, or bug to be fixed. This
  tracking number must be included in the CVS revision history of any
  software change. A set of custom scripts files these revision
  histories to a per-tracking number, easily-accessible file, that can
  be consulted when needed. It is therefore possible to see at a
  glance all the changes related to the corresponding tracking number.
\end{rationale}

% LocalWords:  AdaCore GNAT's GPS IDE GNATUG pragmas VCG GCC standalone CVS
% LocalWords:  AdaCore's
